The training objective for small-lesion segmentation has to satisfy two
requirements (\cref{sec:problem}): a lesion's contribution should not
scale with its size (balance), and a missed lesion should be penalized at
the lesion level (detection). This section reviews prior work according
to the extent to which it meets these requirements.

\paragraph{Voxel- and class-level reweighting.}
Early approaches primarily address the imbalance between foreground and
background. Generalised Dice loss weights each class by the inverse square
of its volume~\cite{sudre_generalised_2017}, whereas Tversky loss assigns a
greater penalty to false negatives than to false
positives~\cite{salehi_tversky_2017}. Focal loss down-weights easy
voxels~\cite{lin_focal_2017}, Focal Tversky loss combines the two
approaches~\cite{abraham_novel_2019}, and DSC++ reduces the overconfidence
associated with Dice loss~\cite{yeung_calibrating_2023}. Region-wise losses
also vary the penalty spatially, either across image
subregions~\cite{chen_adaptive_2023} or through rectified region-wise
maps~\cite{valverde_region-wise_2023}. Large-scale comparisons have found
compound losses to be the most robust choice~\cite{ma_loss_2021,azad_loss_2023}.
However, these methods weight voxels according to their class, difficulty,
or position rather than the lesion to which they belong. Consequently, a
lesion's influence on training remains proportional to its size
(\cref{eq:sl_voxel_loss}), and neither requirement is satisfied.

\paragraph{Instance-level balancing.}
A second group of methods explicitly represents lesions as connected
components. Inverse weighting~\cite{shirokikh_universal_2020} assigns each
voxel a weight inversely proportional to the volume of its lesion. When
applied to Dice, Focal, and asymmetric similarity losses, this approach
improved the detection of liver tumors, lung nodules, and brain metastases,
although delineation quality decreased on the liver dataset. Blob
loss~\cite{kofler_blob_2023}, by contrast, evaluates a base loss separately
for each lesion after masking all other lesions from the domain, and then
averages the resulting losses. With Dice as the base loss, it improved the
lesion-wise F1 score for MS lesions by 5\% relative to soft Dice loss.
CC-DiceCE~\cite{bouteille_learning_2026} builds the same kind of
per-lesion loss on the component-wise evaluation protocol of
CC-Metrics~\cite{jaus_every_2025} and achieves higher recall than blob loss,
albeit with a dataset-dependent reduction in precision. Regularizing both
predictions and labels has also been proposed as a means of addressing
instance imbalance~\cite{aqib_javed_novel_2025}. These methods satisfy the
balance requirement, and the component-adaptive term introduced in this
work follows the same inverse-weighting principle. Nevertheless, they
remain overlap-based: each lesion is assessed according to the quality of
its delineation, coupling detection and delineation within a single term.
None of these methods considers only whether a lesion has been detected.

\paragraph{Lesion-level detection and multiple instance learning.}
Few loss functions directly supervise detection. ALL-Net~\cite{zhang_all-net_2021}
introduces a lesion-wise loss that models MS lesions as equal-sized spheres,
whereas ICI loss~\cite{rachmadi_improving_2024} adds a center-of-instance
term that represents each instance as a fixed-size cube centered at its
center of mass. Both methods remove lesion size from the detection signal,
but rely on surrogate shapes that the prediction must still match.
MIL~\cite{dietterich_solving_1997,ilse_attention-based_2018} expresses the
detection requirement more directly: a bag is considered positive if at
least one of its instances is positive, just as a lesion can be considered
detected if at least one of its voxels is predicted as foreground with high
confidence. In segmentation, however, MIL has been used almost exclusively
to compensate for missing voxel-level labels: pixel scores are pooled to
match image-level labels~\cite{pinheiro_image-level_2015} or bounding-box
constraints~\cite{hsu_weakly_2019}, while recent medical applications have
similarly focused on weakly supervised
segmentation~\cite{liu_object_2026,huang_geometric_2026,pal_deep_2021}. To
the best of our knowledge, MIL has not been used as a per-lesion existence
constraint when full voxel-level annotations are available.

\paragraph{MS lesion segmentation.}
Recent advances in MS lesion segmentation have been driven primarily by
architectural design, ensembling, and training data, while retaining
voxel-level objectives. Examples include a 3D FC-DenseNet trained with an
asymmetric similarity loss~\cite{hashemi_asymmetric_2019}, LST-AI, an
ensemble of three 3D U-Nets trained with Tversky and cross-entropy
loss~\cite{wiltgen_lst-ai_2024}, and the nnU-Net-based
FLAMeS~\cite{dereskewicz_flames_2025}. Although these methods generalize
well across centers, small lesions remain a major source of error: most
lesions missed by FLAMeS were smaller than
10~mm$^3$~\cite{dereskewicz_flames_2025}.

\paragraph{Remaining gap.}
This review highlights three limitations. First, voxel-, class-, and
region-level losses do not explicitly operate on individual lesions.
Second, instance-level losses balance the contributions of lesions of
different sizes but still evaluate each lesion through overlap. They
therefore do not distinguish between detecting a lesion and delineating it
accurately. Third, although MIL directly expresses this detection
criterion, its use in segmentation has largely been confined to weakly
supervised settings. No existing objective combines size-balanced overlap
with an explicit, overlap-independent detection term for each lesion.
\Cref{sec:method} introduces such an objective by adding a
component-adaptive Tversky term for balance and a lesion-level MIL term for
detection to the standard nnU-Net loss~\cite{isensee_nnu-net_2021},
without modifying the architecture.
